\pdfoutput=1
\documentclass[12pt]{article}

\usepackage{amsmath, amssymb, amsthm}
\usepackage{graphicx}
\usepackage{booktabs}
\usepackage{array}
\usepackage{longtable}
\usepackage{calc}
\usepackage{etoolbox}
\makeatletter
\patchcmd\longtable{\par}{\if@noskipsec\mbox{}\fi\par}{}{}
\makeatother
\IfFileExists{footnotehyper.sty}{\usepackage{footnotehyper}}{\usepackage{footnote}}
\makesavenoteenv{longtable}
\usepackage{xurl}
\usepackage[colorlinks=true, linkcolor=blue, citecolor=blue, urlcolor=blue]{hyperref}
\usepackage{geometry}
\geometry{margin=1in}

\makeatletter
\newsavebox\pandoc@box
\newcommand*\pandocbounded[1]{%
  \sbox\pandoc@box{#1}%
  \Gscale@div\@tempa{\textheight}{\dimexpr\ht\pandoc@box+\dp\pandoc@box\relax}%
  \Gscale@div\@tempb{\linewidth}{\wd\pandoc@box}%
  \ifdim\@tempb\p@<\@tempa\p@\let\@tempa\@tempb\fi
  \ifdim\@tempa\p@<\p@\scalebox{\@tempa}{\usebox\pandoc@box}%
  \else\usebox{\pandoc@box}%
  \fi
}
\makeatother
\providecommand{\tightlist}{%
  \setlength{\itemsep}{0pt}\setlength{\parskip}{0pt}}

\title{Tholonic N-D-C Signal Architecture for Accumulation-Biased Grid Trading:\\
\large Design, Implementation, and 59-Day Empirical Results on SOL/USD}
\author{J. W. Milton\\[4pt]{\small Clarity Coalition}}
\date{\small 29 June 2026 \\ v1.0}

\begin{document}
\maketitle

\noindent\textbf{Keywords:} tholonic model; N-D-C framework; grid trading; accumulation strategy; golden ratio; energy cost; imbalance signal; SOL/USD; cryptocurrency; phi spacing

\begin{center}
\textbf{\large CONFIDENTIAL AND PRIVATE. DO NOT DISTRIBUTE.}
\end{center}

\bigskip

\subsection*{Results at a Glance}

\textbf{Period:} 59 days (April 30 to June 29, 2026) \textbar{}
\textbf{Instrument:} SOL/USD \textbar{} \textbf{Runs:} 5,518

\begin{longtable}[]{@{}
  >{\raggedright\arraybackslash}p{(\linewidth - 8\tabcolsep) * \real{0.2000}}
  >{\raggedright\arraybackslash}p{(\linewidth - 8\tabcolsep) * \real{0.2000}}
  >{\raggedright\arraybackslash}p{(\linewidth - 8\tabcolsep) * \real{0.2000}}
  >{\raggedright\arraybackslash}p{(\linewidth - 8\tabcolsep) * \real{0.2000}}
  >{\raggedright\arraybackslash}p{(\linewidth - 8\tabcolsep) * \real{0.2000}}@{}}
\toprule\noalign{}
\begin{minipage}[b]{\linewidth}\raggedright
Metric
\end{minipage} & \begin{minipage}[b]{\linewidth}\raggedright
Start
\end{minipage} & \begin{minipage}[b]{\linewidth}\raggedright
End
\end{minipage} & \begin{minipage}[b]{\linewidth}\raggedright
Change
\end{minipage} & \begin{minipage}[b]{\linewidth}\raggedright
Annualised CAGR
\end{minipage} \\
\midrule\noalign{}
\endhead
\bottomrule\noalign{}
\endlastfoot
SOL held & 556 & 1,279 & +722 (+129.9\%) & \textbf{+17,141\%/yr} (CAGR);
\textasciitilde804\%/yr (simple linear) \\
Total equity (USD) & \$97,179 & \$96,921 & -\$258 (-0.27\%) &
-1.6\%/yr \\
SOL/USD price & \$83.06 & \$71.24 & -\$11.82 (-14.2\%) & -61.3\%/yr \\
Cash balance & \$50,984 & \$5,108 & -\$45,876 deployed & \\
\end{longtable}

\textbf{Acquisition quality:} 722 SOL acquired at an average cost of
\$63.50, against a period low of \$60.48 (3.3\% above the floor) and a
period high of \$98.12. Unrealised gain on acquired position at close:
+\$5,586.

\textbf{Primary objective met.} The bot's goal was to accumulate the
maximum quantity of SOL, not to maximise dollar equity. SOL holdings
more than doubled (+129.9\%) while dollar equity was preserved within
0.3\%, against a backdrop of a 14.2\% price decline in the underlying
asset. The annualised SOL accumulation rate of +17,141\%/yr (compound,
CAGR) or \textasciitilde804\%/yr (simple linear,
\(129.9\% \times 365/59\)) reflects a concentrated 59-day accumulation
phase in a declining market and should not be extrapolated as a
steady-state expectation; it is a measure of execution quality during
the observed period.

\bigskip


\begin{abstract}
This paper describes the design and empirical evaluation of a
limit-order grid trading bot constructed entirely from tholonic N-D-C
principles. The system maps cryptocurrency market state onto the
tholonic triadic structure: bearish constraints are assigned to the D
(Definition) role, bullish contributions to the C (Contribution) role,
and the emergent market coherence to the N (Negotiation) role. Two
derived scalar signals, D-C Imbalance and Energy Cost, govern all
operational decisions including order mix, grid spacing, position
sizing, and regime suppression. Grid level distances are set by a
phi-spaced hierarchy rooted in the natural doubling rate \(\ln(2)\), the
most fundamental tholonic quantum. The bot ran live on SOL/USD for 59
days (April 30 to June 29, 2026), during which SOL fell 14.24\% from
\$83.06 to \$71.24. Over the same period, total equity declined only
0.27\%, and SOL holdings grew from 556 to 1,279 units, a 129.9\%
increase achieved at an average acquisition cost of \$63.50 per SOL,
near the period low of \$60.48. This paper presents the tholonic
derivation of each design component and evaluates the empirical record
against the bot's stated objective: accumulate the maximum quantity of
the target asset.

A central principle of the design is that the asset itself (SOL) is the
correct unit of account, not the dollar. A unit of SOL is structurally
stable: its definitional properties are fixed by the network protocol
and do not change when the market assigns it a different price. The US
dollar, by contrast, is not a stable tholonic unit; its purchasing power
is subject to exogenous forces unrelated to the trading strategy.
Measuring performance in asset units removes this ambiguity and allows
the N-D-C parameters to be defined against a consistent, structurally
grounded substrate.

This paper does not present backtested performance projections,
financial advice, or claims of generalisability beyond the instrument
and period described. All claims are grounded in the event log produced
by the live system.
\end{abstract}

\section{Introduction}\label{introduction}

\textbf{What this paper provides.} A complete account of how tholonic
N-D-C concepts are operationalised as quantitative trading signals; a
derivation of the phi-spaced grid level hierarchy from first principles;
a definition of the Energy and Imbalance signals and their roles in
regime control; and an empirical evaluation of 59 days of live trading
on SOL/USD against an accumulation objective.

\textbf{What this paper does not provide.} Generalised performance
claims, financial return optimisation, price prediction, or any analysis
of the value chain layer of the asset. The bot operates entirely in the
physical-flow layer of the tholonic hierarchy: it measures what the
market is doing (volume, momentum, distance from structural levels) and
responds structurally. It does not forecast price.

\textbf{Organisation.} Section 2 reviews the tholonic N-D-C framework as
applied to market systems; Section 2.1 establishes why the asset unit
rather than the dollar is the correct measurement substrate. Section 3
defines the D and C primitives. Section 4 derives and explains the
Imbalance and Energy signals in depth, including their appearance on the
operational chart. Section 5 derives the phi-spaced grid level
architecture. Section 6 describes the five operational modes. Section 7
covers the auxiliary phi-extension and lambda signals. Section 8
presents the 59-day empirical record. Section 9 discusses the results
against the accumulation objective. Section 10 concludes.

\begin{center}\rule{0.5\linewidth}{0.5pt}\end{center}

\section{The Tholonic N-D-C Framework Applied to Markets}\label{the-tholonic-n-d-c-framework-applied-to-markets}

The tholonic model, developed across
\href{https://github.com/baardev/truevalue/blob/main/docnav/Research/papers/1_recursive-tholonic-five-constants/1_recursive-tholonic-five-constants.pdf}{paper
1} and
\href{https://github.com/baardev/truevalue/blob/main/docnav/Research/papers/3_minimal-recursive-triadic-framework/3_minimal-recursive-triadic-framework.pdf}{paper
3} of this series, holds that any stable system can be decomposed into
three roles:

\begin{itemize}
\tightlist
\item
  \textbf{N (Negotiation):} the emergent, coherent state of the system.
  N is not directly observable; it arises from the balanced interaction
  of D and C.
\item
  \textbf{D (Definition):} the constraints, limitations, and boundaries
  that define what the system is. D is internally focused; it resists
  change and enforces identity.
\item
  \textbf{C (Contribution):} the outputs, flows, and integrations that
  define what the system does. C is externally focused; it expresses
  energy outward.
\end{itemize}

The core sustainability principle of the tholonic model is that a system
is most stable and efficient when \(D \approx C\). When either pole
dominates, the system incurs an energy cost proportional to the square
of the imbalance. This cost is not metaphorical: in physical systems it
corresponds to real dissipation; in market systems it corresponds to
unsustainable overextension in either direction.

\begin{figure}
\centering
\pandocbounded{\includegraphics[width=\linewidth,keepaspectratio]{figures/20_ndc-market-triangle.png}}
\caption{Tholonic N-D-C role assignment for market state. N (blue, top)
is the emergent coherent state. C (red, lower-left) captures bullish
contributions. D (green, lower-right) captures bearish constraints.}
\end{figure}

A market at any given moment can be mapped onto this structure. The
bullish forces that push price upward and volume outward are C-type:
they contribute energy to the system. The bearish forces that resist
further advance and define structural ceilings are D-type: they
constrain and bound. The actual price level, together with the coherence
of that level relative to both poles, is N. When C and D are balanced,
price is at an equilibrium that the system can sustain. When one pole
dominates, the system is overextended and the imbalance creates pressure
toward reversion.

\subsection{Why the Asset Is the Correct Unit of Account}\label{why-the-asset-is-the-correct-unit-of-account}

A foundational choice in this system is to measure success in units of
the asset itself (SOL) rather than in dollars. This choice is not merely
a preference; it follows directly from tholonic principles about what
constitutes a stable, well-defined unit for measurement.

A unit of SOL is a fixed, self-consistent quantity. One SOL remains one
SOL regardless of what price the market assigns to it at any given
moment. Its physical properties as a digital asset, the computation it
represents, the network stake it embodies, the protocol rules that
define it: none of these change when the dollar price moves from \$60 to
\$90 and back. The SOL token is, in tholonic terms, a stable N-state:
its D (the protocol's definitional constraints) and C (its contribution
to the network's operation) are structurally fixed by the protocol
itself. The market price is an external opinion layered on top of that
structure. It is not the thing itself.

The US dollar, by contrast, is not a fixed unit in any tholonic sense.
Its purchasing power shifts continuously with monetary policy,
inflation, and market sentiment. A dollar-denominated performance
measure conflates two separate phenomena: the behavior of the trading
strategy and the behavior of the dollar's own purchasing power. Using
dollars as the unit of success introduces a moving baseline that can
mask or distort the strategy's actual contribution.

This distinction has a direct technical consequence for N-D-C
parameterisation. Because the D and C primitives measure structural
market properties (distance to resistance, volume ratio, trend strength,
support proximity), they can be defined precisely and stably when the
objective unit is the asset. The question ``how much of the asset did we
acquire, and at what structural cost?'' has a clean tholonic answer:
each SOL acquired represents a negotiated outcome between the bullish
forces that wanted to hold it above the buy level and the constraints
that prevented it from rising further. That negotiation is exactly what
the D-C Imbalance signal models.

If success were measured in dollars, the same structural event (a buy
order filling at a phi-spaced level below spot) would have ambiguous
value: it could look like a gain or a loss depending entirely on what
the dollar subsequently does. The asset-denominated objective removes
this ambiguity. The N-D-C parameters are defined against a stable
substrate, and the results can be evaluated cleanly against the question
the system was actually designed to answer: how efficiently is the bot
converting available capital into units of a structurally stable asset?

\begin{center}\rule{0.5\linewidth}{0.5pt}\end{center}

\section{D and C Primitive Definitions}\label{d-and-c-primitive-definitions}

Each of the four D and four C primitives is a scalar in \([0, 10]\),
computed from the most recent closed 15-minute OHLCV bar window and
clamped to prevent single-metric dominance.

\subsection{D Primitives: Bearish Constraints}\label{d-primitives-bearish-constraints}

\begin{longtable}[]{@{}
  >{\raggedright\arraybackslash}p{(\linewidth - 4\tabcolsep) * \real{0.3333}}
  >{\raggedright\arraybackslash}p{(\linewidth - 4\tabcolsep) * \real{0.3333}}
  >{\raggedright\arraybackslash}p{(\linewidth - 4\tabcolsep) * \real{0.3333}}@{}}
\toprule\noalign{}
\begin{minipage}[b]{\linewidth}\raggedright
Primitive
\end{minipage} & \begin{minipage}[b]{\linewidth}\raggedright
Formula
\end{minipage} & \begin{minipage}[b]{\linewidth}\raggedright
Interpretation
\end{minipage} \\
\midrule\noalign{}
\endhead
\bottomrule\noalign{}
\endlastfoot
Resistance distance &
\(\text{clamp}\!\left(\frac{R - P}{P} \times 100\right)\) & Distance to
nearest swing high above price. High resistance nearby = high D. \\
RSI overbought &
\(\text{clamp}\!\left(\frac{\text{RSI} - 50}{5}\right)\) & RSI above 50
contributes to D; below 50 contributes zero. \\
Realised volatility &
\(\text{clamp}\!\left(\frac{\sigma_r}{0.05}\right)\) & 30-bar rolling
realised vol normalised to a 5\% baseline. High vol = high uncertainty =
high D. \\
Funding rate & \(\text{clamp}\!\left(\max(f, 0) \times 100\right)\) &
Positive funding on perpetuals indicates excess long positioning =
bearish constraint on further advance. \\
\end{longtable}

\subsection{C Primitives: Bullish Contributions}\label{c-primitives-bullish-contributions}

\begin{longtable}[]{@{}
  >{\raggedright\arraybackslash}p{(\linewidth - 4\tabcolsep) * \real{0.3333}}
  >{\raggedright\arraybackslash}p{(\linewidth - 4\tabcolsep) * \real{0.3333}}
  >{\raggedright\arraybackslash}p{(\linewidth - 4\tabcolsep) * \real{0.3333}}@{}}
\toprule\noalign{}
\begin{minipage}[b]{\linewidth}\raggedright
Primitive
\end{minipage} & \begin{minipage}[b]{\linewidth}\raggedright
Formula
\end{minipage} & \begin{minipage}[b]{\linewidth}\raggedright
Interpretation
\end{minipage} \\
\midrule\noalign{}
\endhead
\bottomrule\noalign{}
\endlastfoot
Volume ratio &
\(\text{clamp}\!\left(5 \times \frac{V_t}{\bar{V}_{20}}\right)\) &
Current bar volume relative to 20-bar mean. Rising volume = rising C. \\
Momentum score &
\(\text{clamp}\!\left(\frac{P - \text{EMA}_{50}}{\text{ATR}}\right)\) &
Price above its 50-bar EMA = bullish momentum. Normalised by ATR. \\
Trend strength & \(\text{clamp}\!\left(\frac{\text{ADX}}{10}\right)\) &
ADX normalised to \([0, 10]\). Strong directional trend = high C. \\
Support distance &
\(\text{clamp}\!\left(10 - \frac{P - S}{P} \times 100\right)\) &
Proximity to nearest swing low below price. Close to support = high C
(limited downside). \\
\end{longtable}

\(D_{\text{total}} = \sum D_i\), \(\quad C_{\text{total}} = \sum C_i\)

Both totals have a theoretical range of \([0, 40]\) given four clamped
\([0, 10]\) components.

\begin{center}\rule{0.5\linewidth}{0.5pt}\end{center}

\section{Imbalance and Energy: The Core Signal Pair}\label{imbalance-and-energy-the-core-signal-pair}

These two derived scalars are the engine of every operational decision
in the bot. They appear on the live operational chart as the two lower
panels (labelled ``NDC Imbalance'' and, implicitly, the energy threshold
crossings).

\subsection{Imbalance}\label{imbalance}

\[\text{Imbalance} = \frac{C_{\text{total}} - D_{\text{total}}}{\max(C_{\text{total}}, D_{\text{total}})}\]

Imbalance is a signed, normalised measure of which pole dominates, on a
scale from \(-1\) (fully D-dominant, maximum bearish constraint) to
\(+1\) (fully C-dominant, maximum bullish contribution).

At zero, \(C = D\) and the system is in tholonic balance. The threshold
\(\pm 0.0347\) is not arbitrary: it is \(\ln(2)/20\), half the base
tholonic quantum, used as the minimum meaningful deviation from balance
before the regime shifts.

The interpretation at each regime boundary:

\begin{longtable}[]{@{}
  >{\raggedright\arraybackslash}p{(\linewidth - 4\tabcolsep) * \real{0.3333}}
  >{\raggedright\arraybackslash}p{(\linewidth - 4\tabcolsep) * \real{0.3333}}
  >{\raggedright\arraybackslash}p{(\linewidth - 4\tabcolsep) * \real{0.3333}}@{}}
\toprule\noalign{}
\begin{minipage}[b]{\linewidth}\raggedright
Imbalance
\end{minipage} & \begin{minipage}[b]{\linewidth}\raggedright
Market meaning
\end{minipage} & \begin{minipage}[b]{\linewidth}\raggedright
Bot response
\end{minipage} \\
\midrule\noalign{}
\endhead
\bottomrule\noalign{}
\endlastfoot
\(> +0.0347\) & C overwhelms D; market in bullish overexpression &
Distribute: 3 buys, 5 sells \\
\([-0.0347,\, +0.0347]\) & Near balance; no dominant pole & Neutral: 5
buys, 3 sells (accumulation default) \\
\(< -0.0347\) & D overwhelms C; market compressed by bearish constraint
& Accumulate: 5 buys, 1 sell \\
\end{longtable}

Crucially, the neutral default is asymmetric: 5 buys and 3 sells rather
than equal counts. This encodes the stated objective (accumulate the
asset) at the structural level. Even when the signal is inconclusive,
the bot's default behavior is accumulation-biased.

\subsection{Energy Cost}\label{energy-cost}

\[E = |\text{Imbalance}|^2 \times D_{\text{max}}^2 + E_{\text{base}}\]

In a simplified but operationally equivalent form used in the code:

\[E = \left(\frac{|C - D|}{\max(C, D)}\right)^2 + E_{\text{base}}\]

where \(E_{\text{base}} = 10\).

Energy quantifies the cost of maintaining the current imbalanced state.
It is not volatility; it is the tholonic cost of imbalance: a system far
from \(D = C\) must expend energy to maintain that deviation, and the
cost grows quadratically with the deviation. A balanced system
(\(D = C\)) has only the base cost \(E_{\text{base}}\). A maximally
imbalanced system (one pole at zero) has
\(E = 1 + E_{\text{base}} = 11\) at the formula level, but the quadratic
scaling means large sustained deviations quickly reach the regime
threshold.

The energy threshold \(E_{\text{max}} = 135\) defines the boundary above
which the system is in a high-cost, unstable state. When
\(E > E_{\text{max}}\):

\begin{itemize}
\tightlist
\item
  If Imbalance \(> 0\) (C-dominant spike): the market is in a bull
  overextension. The bot enters \textbf{buy-only mode}, suppressing all
  sell orders and accumulating through the rally.
\item
  If Imbalance \(\leq 0\) (D-dominant spike): the market is in chaos or
  a bear crush. The bot \textbf{pauses} entirely, placing no orders.
\end{itemize}

This asymmetry is again structural: C-dominant high-energy states are
accumulation opportunities; D-dominant high-energy states are danger
zones.

\subsection{The N State}\label{the-n-state}

\[N = \sqrt{D_{\text{total}} \times C_{\text{total}}} \times B\]

where the balance factor \(B = \dfrac{1}{1 + |D - C| / \max(D, C)}\).

N is the geometric mean of the two poles, penalised by their imbalance.
This formulation is consistent with the tholonic principle established
in
\href{https://github.com/baardev/truevalue/blob/main/docnav/Research/papers/1_recursive-tholonic-five-constants/1_recursive-tholonic-five-constants.pdf}{paper
1}: N maximises when \(D = C\), and collapses toward zero as either pole
approaches zero. N is logged at each run but does not directly control
order placement; it is the diagnostic of overall system coherence.

\subsection{Reading the Operational Chart}\label{reading-the-operational-chart}

The screenshot accompanying this paper shows the live operational
display. The lower chart panel, labelled ``NDC Imbalance,'' plots the
Imbalance signal as a time series against its \(\pm 0.0347\) threshold
bands. The green-filled region above zero reflects sustained
C-dominance; the narrow red regions reflect the brief D-dominant
episodes during the May price decline. The flat band near zero in quiet
periods corresponds to the neutral regime.

The energy signal does not appear as a separate visual panel in the
standard chart display, but its effect is visible in the order activity:
the 15 buy-only events and the absence of any full pause events
correspond to the 953 runs (17.3\% of the total) where energy exceeded
\(E_{\text{max}} = 135\).

\begin{figure}
\centering
\pandocbounded{\includegraphics[width=\linewidth,keepaspectratio]{figures/20_energy-imbalance-timeseries.png}}
\caption{SOL/USD price (top), D-C Imbalance (middle), and Energy Cost
(bottom) over the 59-day run. Red bars in the imbalance panel indicate
C-dominant conditions; green bars indicate D-dominant. The horizontal
dashed line in the energy panel marks the pause/hold threshold.}
\end{figure}

\begin{center}\rule{0.5\linewidth}{0.5pt}\end{center}

\section{Grid Architecture: Phi-Spaced Level Hierarchy}\label{grid-architecture-phi-spaced-level-hierarchy}

\subsection{The Base Quantum}\label{the-base-quantum}

The grid level spacings are derived from the tholonic base quantum:

\[u = \frac{\ln 2}{10} \times s\]

where \(s\) is a timeframe scale factor (\(s = 0.35\) for 15-minute
bars; \(s = 1.0\) for 4-hour bars). The use of \(\ln 2\) is not
arbitrary: it is the natural doubling rate, the minimal quantum of
binary distinction in the tholonic framework. At \(s = 0.35\), the base
quantum is \(u \approx 0.00243\), or 0.243\% of the current price per
level.

\subsection{The Five Phi-Spaced Levels}\label{the-five-phi-spaced-levels}

Five grid levels are placed at the following distances from the current
price, in both directions:

\[\left\{\frac{u}{\varphi^2},\; \frac{u}{\varphi},\; u,\; u\varphi,\; u\varphi^2\right\}\]

where \(\varphi = (1 + \sqrt{5})/2 \approx 1.61803\) is the golden
ratio.

At \(s = 0.35\), these evaluate to approximately:

\begin{longtable}[]{@{}lll@{}}
\toprule\noalign{}
Level & Formula & Distance from spot \\
\midrule\noalign{}
\endhead
\bottomrule\noalign{}
\endlastfoot
1 & \(u / \varphi^2\) & 0.93\% \\
2 & \(u / \varphi\) & 1.50\% \\
3 & \(u\) & 2.43\% \\
4 & \(u \cdot \varphi\) & 3.93\% \\
5 & \(u \cdot \varphi^2\) & 6.35\% \\
\end{longtable}

The ratio of any adjacent level pair is \(\varphi\). This makes the grid
self-similar: zooming into the spacing between any two levels reveals
the same proportional structure as the whole. This is the direct market
application of the tholonic principle that stable, sustainable
hierarchical structures are self-similar at all scales.

\begin{figure}
\centering
\pandocbounded{\includegraphics[width=\linewidth,keepaspectratio]{figures/20_phi-grid-levels.png}}
\caption{Phi-spaced grid level diagram showing buy levels (red, below
spot) and sell levels (green, above spot) with their tholonic formula
labels and percentage distances from the current price.}
\end{figure}

\begin{center}\rule{0.5\linewidth}{0.5pt}\end{center}

\section{Operational Modes and Regime Logic}\label{operational-modes-and-regime-logic}

\subsection{Hysteresis}\label{hysteresis}

Regime changes require the signal to exceed the threshold for at least
two consecutive bars before the mode flips. This prevents oscillation
around the threshold boundary from causing excessive order cancellation
and replacement. The hysteresis band effectively smooths the mode signal
while preserving responsiveness to sustained changes.

\subsection{The Five Modes}\label{the-five-modes}

\textbf{Distribute} (\(\text{Imbalance} > +0.0347\), normal energy): 3
buy levels below spot, 5 sell levels above spot. The bot is reading a
C-dominant environment and positions to capture mean reversion: few buys
(limited upside buying interest) and more sells (take profit into
strength). In a sustained downtrend, sell levels placed above a
declining price do not fill; the buy levels fill repeatedly. This
explains the paradox that a bot running 87\% of the time in distribute
mode still accumulated net SOL during a declining market.

\textbf{Neutral} (\(|\text{Imbalance}| \leq 0.0347\)): 5 buy levels, 3
sell levels. The default accumulation-biased grid.

\textbf{Accumulate} (\(\text{Imbalance} < -0.0347\), normal energy): 5
buy levels, 1 sell level. Near-suppression of selling during D-dominant
compression.

\textbf{Buy-only} (\(E > E_{\text{max}}\) and \(\text{Imbalance} > 0\)):
5 buy levels, 0 sell levels. The system reads a high-energy C-dominant
spike as a bull overextension that should be accumulated through, not
sold into.

\textbf{Pause} (\(E > E_{\text{max}}\) and \(\text{Imbalance} \leq 0\)):
0 orders. High-energy D-dominant states are structurally unpredictable;
the bot steps aside entirely.

\begin{center}\rule{0.5\linewidth}{0.5pt}\end{center}

\section{Auxiliary Signals}\label{auxiliary-signals}

\subsection{Phi-Extension Signal}\label{phi-extension-signal}

The ratio of successive swing highs is compared to \(\varphi\):

\[r = \frac{H_{t}}{H_{t-1}}\]

If \(r > \varphi + 0.15\): the market is overextended relative to its
natural growth rate. One buy level is removed and one sell level added
(tilt toward distribution).

If \(r < \varphi - 0.15\): the market is under-extended. One buy level
is added and one sell level removed (tilt toward accumulation).

This signal fired in only 4.4\% of runs during the observed period
(phi\_signal = +1, under-extended), consistent with a sustained
declining market that never produced a blow-off top.

\subsection{Lambda: Volume Recovery Signal}\label{lambda-volume-recovery-signal}

Lambda detects volume recovering from a recent trough:

\[\lambda = \min\!\left(1,\; \frac{V_{\text{ratio}} - 1}{2}\right) \quad \text{when volume is recovering}\]

Buy order sizes are scaled by \((1 + 0.618 \times \lambda)\), where the
weight \(0.618 = 1/\varphi\) is the tholonic balance ratio. Lambda fired
in 1,291 of 5,518 runs (23.4\%), scaling up buy sizes during
volume-recovery events to front-run anticipated C-signal strengthening.

\begin{center}\rule{0.5\linewidth}{0.5pt}\end{center}

\section{Empirical Results: 59-Day SOL/USD Run}\label{empirical-results-59-day-solusd-run}

\subsection{Setup}\label{setup}

The bot ran live on the Alpaca cryptocurrency API, trading SOL/USD on a
15-minute execution cadence. The signal timeframe was 15-minute OHLCV
bars with a 100-bar lookback. A 5-minute reclaim trigger filter required
lower-timeframe confirmation before order placement. Enhanced NDC risk
controls were enabled, including adaptive ATR-based spacing,
inventory-aware sizing, and the no-trade zone suppression.

\subsection{Summary Statistics}\label{summary-statistics}

\begin{longtable}[]{@{}ll@{}}
\toprule\noalign{}
Metric & Value \\
\midrule\noalign{}
\endhead
\bottomrule\noalign{}
\endlastfoot
Period & April 30 to June 29, 2026 (59 days) \\
Total runs & 5,518 (15-minute cadence) \\
SOL/USD price change & \$83.06 to \$71.24 (-14.24\%) \\
SOL price low & \$60.48 \\
Starting equity & \$97,179 \\
Ending equity & \$96,921 (-0.27\%) \\
Equity peak & \$108,194 \\
Equity trough & \$83,635 \\
Starting cash & \$50,984 \\
Ending cash & \$5,108 \\
Starting SOL held & 556 \\
Ending SOL held & 1,279 \\
SOL gained & +722 (+129.9\%) \\
Implied average acquisition price & \$63.50/SOL \\
Unrealised gain on acquired SOL & +\$5,586 \\
Total orders placed & 28,963 \\
Order success rate & 71.5\% \\
Total buy orders & 16,456 \\
Total sell orders & 24,048 \\
\end{longtable}

\subsection{Mode Distribution}\label{mode-distribution}

The distribute mode dominated at 86.8\% of runs, reflecting sustained
C-dominance (mean imbalance +0.44) even as price declined. Accumulate
mode ran for 6.5\% of runs, concentrated during the price troughs in May
and June. Buy-only mode fired 15 times during brief high-energy
C-dominant spikes.

\begin{figure}
\centering
\pandocbounded{\includegraphics[width=\linewidth,keepaspectratio]{figures/20_mode-distribution.png}}
\caption{Operational mode distribution across 5,518 runs (left) and the
distribution of D-C Imbalance values over the full period (right). The
imbalance distribution is strongly right-skewed, confirming persistent
C-dominance.}
\end{figure}

\subsection{Accumulation Record}\label{accumulation-record}

The bot converted \$45,875 in cash into 722 additional SOL at an average
cost of \$63.50 per unit, 23.5\% below the entry price of \$83.06 and
only \$3.02 above the period low of \$60.48. The accumulation rate
accelerated in June as price fell further and lambda-scaling increased
buy sizes at volume-recovery inflection points.

\begin{figure}
\centering
\pandocbounded{\includegraphics[width=\linewidth,keepaspectratio]{figures/20_accumulation-cash-equity.png}}
\caption{SOL held (top, red), cash balance (middle, green), and total
equity (bottom, blue) over the 59-day run. Cash was systematically
deployed into SOL as price declined, with SOL holdings more than
doubling while equity remained nearly flat.}
\end{figure}

\begin{center}\rule{0.5\linewidth}{0.5pt}\end{center}

\section{Discussion}\label{discussion}

\subsection{Accumulation Objective Assessment}\label{accumulation-objective-assessment}

The stated objective was to accumulate the maximum quantity of SOL
rather than to maximise dollar-denominated equity. By that measure, the
result is strong: 129.9\% growth in SOL holdings over a period where
holding SOL in cash-equivalent terms would have cost 14.24\% in dollar
value. The bot converted that paper loss into asset accumulation,
acquiring 722 SOL at below-market cost.

The near-flat equity (-0.27\%) is not the primary measure of success but
confirms that the accumulation was achieved without significant
dollar-denominated destruction. The equity low of \$83,635 (in late May)
coincided with the price trough, at which point the bot's unrealised
losses on earlier SOL purchases were at their maximum. Recovery in June
reflected both price stabilisation and the accumulation of lower-cost
SOL reducing the average basis.

\subsection{The Distribute-Mode Accumulation Paradox}\label{the-distribute-mode-accumulation-paradox}

The dominant finding is that the bot ran 87\% of the time in distribute
mode while net-accumulating SOL. This is explained by the asymmetric
fill dynamics of a limit-order grid in a trending market. In distribute
mode, sell orders are placed above the current price. When price is
declining, those sells are never reached. The buy orders placed below
spot, however, are reached on every dip. The net effect is that
distribute mode in a declining market becomes a systematic buy program,
despite its label. The tholonic framework did not fail here; the label
``distribute'' refers to the bot's intent given the signal state, not to
the actual fills achieved.

A recalibration of the D primitive weights would correct the signal: if
the D parameters better captured the declining trend (for example, by
weighting MACD histogram or the rate of change of resistance levels),
imbalance would have shifted toward D-dominance or neutral more often,
reducing the distribute-mode dominance. This is a parameter tuning
question, not a structural flaw.

\subsection{Cash Exhaustion Risk}\label{cash-exhaustion-risk}

Cash declined from \$50,984 to \$5,108 over the period. With cash nearly
depleted, the bot's capacity to place buy orders is now severely
constrained. If price continues to decline, the accumulation rate will
fall to whatever sell fills can be captured at higher prices and
recycled. The \texttt{max\_deploy\_usd} and
\texttt{capital\_per\_level\_usd} parameters should be reviewed and
tightened to preserve a minimum cash reserve.

\subsection{Order Rejection Rate}\label{order-rejection-rate}

The 28.5\% order rejection rate is elevated and represents wasted API
calls and potential missed fills. The likely cause is minimum notional
violations: as SOL price declined, the fixed dollar-per-level size
(\$1,794 default) converted to increasingly large SOL quantities, and
the resulting sell orders at higher price levels sometimes fell below
the exchange's minimum notional. Dynamic sizing proportional to current
price, rather than fixed dollar amounts, would reduce rejections.

\begin{center}\rule{0.5\linewidth}{0.5pt}\end{center}

\section{Conclusion}\label{conclusion}

The tholonic N-D-C framework provides a structurally coherent basis for
market signal design. Mapping bearish constraints to D, bullish
contributions to C, and market coherence to N produces an operationally
interpretable signal pair: Imbalance measures which pole dominates;
Energy measures the cost of that domination. Together they govern all
regime transitions in a principled, derivable way rather than through
empirically tuned moving-average crossovers.

The phi-spaced grid level hierarchy, derived from \(\ln(2)\) and
\(\varphi\), places limit orders at self-similar distances that
correspond to the natural structure of a recursive triadic system. The
golden ratio appears not as decoration but as the structural attractor
of the level spacing, directly analogous to its appearance in the
convergence hierarchies documented in
\href{https://github.com/baardev/truevalue/blob/main/docnav/Research/papers/1_recursive-tholonic-five-constants/1_recursive-tholonic-five-constants.pdf}{paper
1} and
\href{https://github.com/baardev/truevalue/blob/main/docnav/Research/papers/11_tholonic-seed-space-power-of-two-hierarchy/11_tholonic-seed-space-power-of-two-hierarchy.pdf}{paper
11} of this series.

The 59-day live record on SOL/USD confirms that the system achieved its
accumulation objective: 129.9\% growth in SOL holdings at an average
cost near the period low, with equity preservation of -0.27\% against a
-14.24\% spot price decline. The open questions (D-parameter
recalibration, cash reserve management, order sizing dynamics) are
parameter-level refinements, not structural concerns.

\begin{center}\rule{0.5\linewidth}{0.5pt}\end{center}

\section{References}\label{references}

{[}Mil26a{]} Milton, J. W. \emph{Emergence of Classical Constants from a
Minimal Recursive Triadic Framework.} Clarity Coalition, paper 1 in this
series, 2026.

\url{https://github.com/baardev/truevalue/blob/main/docnav/Research/papers/1_recursive-tholonic-five-constants/1_recursive-tholonic-five-constants.pdf}

{[}Mil26b{]} Milton, J. W. \emph{A Minimal Recursive Triadic Framework
for Self-Similar Hierarchical Systems.} Clarity Coalition, paper 3 in
this series, 2026.

\url{https://github.com/baardev/truevalue/blob/main/docnav/Research/papers/3_minimal-recursive-triadic-framework/3_minimal-recursive-triadic-framework.pdf}

{[}Mil26c{]} Milton, J. W. \emph{Power-of-Two Convergence Counts in the
Tholonic Seed Space: A Complete Hierarchy of Classical Constants.}
Clarity Coalition, paper 11 in this series, 2026.

\url{https://github.com/baardev/truevalue/blob/main/docnav/Research/papers/11_tholonic-seed-space-power-of-two-hierarchy/11_tholonic-seed-space-power-of-two-hierarchy.pdf}

\end{document}
